\documentclass{amsart}
% Adapted from the supplied FirstVersion.tex (Ryuya Hora template).
% The environment-specific dmjhira font declaration is omitted for portability.
\usepackage[left=2cm,right=2cm,top=2.2cm,bottom=2.2cm]{geometry}
\usepackage[utf8]{inputenc}
\usepackage[T1]{fontenc}
\usepackage{amsfonts,amsthm,amssymb,mathtools,etoolbox,stmaryrd}
\usepackage[colorlinks=true,urlcolor=blue,linkcolor=blue,citecolor=blue]{hyperref}
\usepackage{tikz,tikz-cd}
\usetikzlibrary{calc,arrows.meta,positioning,fit,backgrounds}
\usepackage{cleveref}
\usepackage{xcolor}
\usepackage{array,booktabs,longtable,tabularx}
\usepackage{enumitem}
\usepackage{microtype}
\usepackage{xurl}
\usepackage{listings}
\usepackage[style=alphabetic,sorting=nyt,maxnames=5]{biblatex}
\renewbibmacro{in:}{}
\addbibresource{CommonBiblio20240922.bib}
\addbibresource{gsh_additions.bib}

\tikzset{pullback/.style={minimum size=1.2ex,path picture={
\draw[opacity=1,black,-,#1] (-0.5ex,-0.5ex) -- (0.5ex,-0.5ex) -- (0.5ex,0.5ex);%
}}}

\theoremstyle{plain}
\newtheorem{theorem}{Theorem}[subsection]
\newtheorem{proposition}[theorem]{Proposition}
\newtheorem{lemma}[theorem]{Lemma}
\newtheorem{corollary}[theorem]{Corollary}
\newtheorem{todo}[theorem]{Todo}
\newtheorem{conjecture}[theorem]{Conjecture}
\newtheorem{fact}[theorem]{Fact}
\newtheorem{claim}{Claim}[theorem]
\crefname{claim}{Claim}{Claims}
\renewcommand{\theclaim}{\thetheorem.\alph{claim}}

\theoremstyle{definition}
\newtheorem{example}[theorem]{Example}
\newtheorem{definition}[theorem]{Definition}
\newtheorem{notation}[theorem]{Notation}
\newtheorem{puzzle}[theorem]{Puzzle}
\newtheorem{idea}[theorem]{Idea}
\newtheorem{question}[theorem]{Question}
\newtheorem{exercise}[theorem]{Exercise}
\newtheorem{protocol}[theorem]{Protocol}

\theoremstyle{remark}
\newtheorem{remark}[theorem]{Remark}
\newtheorem{warning}[theorem]{Warning}

\newcommand{\dq}[1]{``#1''}
\newcommand{\memo}[1]{\textcolor{red}{memo: #1}}
\newcommand{\invmemo}[1]{}
\newcommand{\horamemo}[1]{\textcolor{green!60!black}{hora: #1}}
\newcommand{\para}[1]{\paragraph{\textbf{#1}}}
\newcommand{\N}{\mathbb{N}}
\newcommand{\Z}{\mathbb{Z}}
\newcommand{\Q}{\mathbb{Q}}
\newcommand{\R}{\mathbb{R}}
\newcommand{\C}{\mathcal{C}}
\newcommand{\D}{\mathcal{D}}
\newcommand{\E}{\mathcal{E}}
\newcommand{\F}{\mathbf{F}}
\renewcommand{\L}{\mathcal{L}}
\newcommand{\id}{\mathrm{id}}
\newcommand{\op}{\mathrm{op}}
\newcommand{\ob}{\mathrm{ob}}
\newcommand{\Set}{\mathbf{Set}}
\newcommand{\FinSet}{\mathbf{FinSet}}
\newcommand{\Func}[2]{[#1,#2]}
\newcommand{\abs}[1]{\left|#1\right|}
\newcommand{\demph}[1]{\textit{#1}}
\newcommand{\Pow}{\mathcal{P}}
\newcommand{\Words}[1]{#1^{\ast}}
\newcommand{\eps}{\varepsilon}
\newcommand{\sh}{\operatorname{sh}}
\newcommand{\Syn}{\operatorname{Syn}}
\newcommand{\Aper}{\mathbf{A}}
\newcommand{\Afour}{A_{4}}
\newcommand{\Afive}{A_{5}}
\newcommand{\Dicthree}{\operatorname{Dic}_{3}}
\newcommand{\heightone}{\mathsf{H}_{1}}
\newcommand{\code}[1]{\texttt{\detokenize{#1}}}
\newcommand{\status}[1]{\textsf{#1}}

\definecolor{codegray}{gray}{0.97}
\lstdefinestyle{workshop}{
  basicstyle=\ttfamily\small,
  backgroundcolor=\color{codegray},
  frame=single,
  framerule=0.2pt,
  columns=fullflexible,
  keepspaces=true,
  breaklines=true,
  showstringspaces=false,
  upquote=true,
  literate={α}{{$\alpha$}}1
}
\lstset{style=workshop}

\setlist[itemize]{topsep=3pt,itemsep=2pt,parsep=0pt}
\setlist[enumerate]{topsep=3pt,itemsep=2pt,parsep=0pt}
\setcounter{tocdepth}{2}
\setcounter{secnumdepth}{3}


\title{Finite Groups and Cohomology for Formal-Language Theorists}
\author{Ryuya Hora and the ZMC Generalized Star-Height Working Group}
\address{A working primer for the ZMC conference, ZEN University}
\date{22 July 2026}
\subjclass[2020]{20D10, 20J06, 68Q45, 20M35}
\keywords{finite groups, semidirect products, group cohomology, group languages, generalized star height, $A_5$}

\begin{document}
\begin{abstract}
This primer supplies the finite-group, extension, representation, and low-dimensional cohomology background needed by formal-language researchers working on generalized star height.  The emphasis is operational: how group decompositions affect evaluation of words, why class-two and $A\rtimes C_2$ formulas lead to counting languages, where the $C_3$ argument fails, and what concrete data would make a cohomological route meaningful.  The groups $A_4$, $\operatorname{Dic}_3$, and $A_5$ are treated as successive research test cases.  Every chapter ends with translation tasks designed to produce exact language statements or machine-checkable artifacts.
\end{abstract}
\maketitle
\tableofcontents

\section{Why finite-group structure matters}

A language recognized by a finite group $G$ has the form
\[
  L=\varphi^{-1}(P),\qquad \varphi:A^\ast\to G,
\]
where $\varphi$ is determined by the images of letters and $P\subseteq G$.  If the multiplication in $G$ can be expressed through finitely many modular statistics of a word, then one may hope to build a generalized star-height-one expression for each accepting fiber.

The successful literature does precisely this for several group families \cite{pinStraubingTherien1992}.  Group theory supplies normal forms and collection formulas; formal-language theory must then express the resulting statistics using Boolean operations, concatenation, complement, and a single nonnested star.  The second step is not automatic.

\begin{warning}
A structural decomposition $G=N\rtimes Q$ is not by itself a height-one theorem.  It must be accompanied by an unambiguous or otherwise controlled map from words to the $N$- and $Q$-data in the multiplication formula.
\end{warning}

\section{Finite-group essentials in recognition form}

\subsection{Generated image and accepting subsets}

A morphism $\varphi:A^\ast\to G$ need not be surjective.  Its image is the subgroup generated by the letter images.  Thus proofs about all $G$-recognized languages may often reduce to all subgroups of $G$, but the direction must be stated carefully.  The accepting subset $P$ is arbitrary; a proof only for the identity fiber or for conjugacy-invariant $P$ is insufficient unless a closure reduction handles all subsets.

Left or right translating $P$ corresponds to quotienting a language by a word when the translating element is represented by a word.  Automorphisms of $G$ may identify recognition data, but a computational symmetry reduction must prove that expression height is preserved under the induced relabeling.

\subsection{Normal subgroups and extensions}

For $N\trianglelefteq G$, the quotient map gives
\[
1\longrightarrow N\longrightarrow G\stackrel{\pi}{\longrightarrow}Q\longrightarrow1.
\]
Given a word $w$, the quotient component $\pi(\varphi(w))$ is easy to evaluate.  Recovering the $N$ component requires choices of coset representatives.  A section $s:Q\to G$ yields a factor set
\[
  f(q_1,q_2)=s(q_1)s(q_2)s(q_1q_2)^{-1}\in N.
\]
Associativity imposes a cocycle identity, twisted by the conjugation action.  This elementary formula is the most concrete point at which cohomology could enter the star-height problem.

\subsection{Split extensions and semidirect products}

If the section is a homomorphism, the extension splits and
\[
G\cong N\rtimes_\theta Q,
\]
where $\theta:Q\to\operatorname{Aut}(N)$.  Write elements as $(n,q)$, with
\[
(n,q)(n',q')=(n\,\theta(q)(n'),qq').
\]
For a word with letter images $(n_a,q_a)$, the quotient product is $q_{a_1}\cdots q_{a_k}$ and the normal component is
\[
  n_{a_1}\,	heta(q_{a_1})(n_{a_2})\,	heta(q_{a_1}q_{a_2})(n_{a_3})\cdots.
\]
This is a weighted path statistic in the quotient automaton.  A language proof must express it uniformly.

\section{Commutative and class-two groups}

\subsection{Abelian groups}

For an abelian group, word evaluation is independent of order:
\[
  \varphi(w)=\prod_{a\in A}\varphi(a)^{\#_a(w)}.
\]
Because $G$ is finite, only the letter counts modulo suitable exponents matter.  This reduces recognition to a Boolean combination of modular count conditions.  Pin--Straubing--Th\'erien prove the resulting languages have generalized star-height at most one \cite{pinStraubingTherien1992}.

The formal-language burden is to express modular counting with one layer of star.  Complement is crucial: languages such as words containing no $a$ are star-free because
\[
  \{w:\#_a(w)=0\}=\bigl(A^\ast a A^\ast\bigr)^c,
  \qquad A^\ast=\varnothing^c.
\]

\subsection{Nilpotency class two}

A group has nilpotency class at most two when $[G,G]\subseteq Z(G)$.  Collection then stops after commutator corrections.  Fix an ordered alphabet.  Reordering a word into grouped letters gives
\[
\varphi(w)=
  \prod_a g_a^{\#_a(w)}
  \prod_{a<b}[g_b,g_a]^{I_{a,b}(w)},
\]
where $I_{a,b}(w)$ counts inversions of type $(b,a)$, up to the adopted commutator convention.  Since commutators are central and finite-order, only modular values matter.

This formula illustrates a reusable pattern:
\begin{enumerate}
  \item prove a finite algebraic normal form;
  \item identify exact path/subword statistics;
  \item show each statistic language has a height-one expression;
  \item combine fibers by Boolean operations.
\end{enumerate}
The third step is where formal-language expertise is indispensable.

\begin{exercise}
For the Heisenberg group over $\Z/m\Z$, choose letter images and derive the matrix product in terms of letter counts and ordered-pair counts.  State the exact modular languages whose intersection gives the identity fiber.
\end{exercise}

\section{Solvability, nilpotence, and why $A_5$ is different}

\subsection{Derived and lower central series}

The derived series is
\[
  G^{(0)}=G,\qquad G^{(i+1)}=[G^{(i)},G^{(i)}].
\]
A group is solvable if this reaches the identity.  The lower central series is
\[
  \gamma_1(G)=G,\qquad \gamma_{i+1}(G)=[\gamma_i(G),G],
\]
and $G$ is nilpotent if this reaches the identity.  Nilpotent groups are solvable, but not conversely.

Induction along an abelian normal series is tempting for star height.  It is not currently a general theorem: every extension step needs a language-theoretic transport mechanism.  The first missing cyclic quotient already appears at order three.

\subsection{Simple groups}

A nontrivial group is simple when it has no nontrivial proper normal subgroup.  $A_5$ is the smallest nonabelian simple finite group.  Extension induction cannot break it into smaller normal layers.  This makes it a clean test of whether height-one phenomena are fundamentally solvable-group phenomena or whether subgroup/coset and representation data can replace normal-series induction.

\section{The order-twelve bottleneck}

\subsection{$A_4=V_4\rtimes C_3$}

The normal Klein four subgroup is
\[
  V_4=\{1,(12)(34),(13)(24),(14)(23)\}.
\]
A $3$-cycle conjugates its three nontrivial elements cyclically.  Thus
\[
  A_4\cong V_4\rtimes C_3.
\]
This group is solvable and small, but it lies outside the published $C_2$-quotient mechanism that Bourne attempted to generalize \cite{bourne2017counting}.

For a recognition morphism $\varphi:A^\ast\to A_4$, projection to $C_3$ gives a three-state quotient automaton.  The $V_4$ component is a sum of letter contributions twisted by the prefix quotient state.  It is therefore an arrow-count statistic on this cyclic automaton.  The failed argument sought a unique decomposition into designated cycles; overlapping occurrences destroy uniqueness.

\subsection{$\operatorname{Dic}_3=C_3\rtimes C_4$}

Let $a$ have order three and $b$ order four, with $b^{-1}ab=a^{-1}$.  Every element has a normal form $a^i b^j$, and
\[
  (a^i b^j)(a^{i'}b^{j'})
  =a^{i+(-1)^j i'}b^{j+j'}.
\]
The action factors through $C_4\twoheadrightarrow C_2$, but the quotient automaton has four states.  This group tests whether one can exploit the action kernel while retaining the quotient path information needed for acceptance.

\begin{exercise}
Derive the normal-component formula for a word in $C_3\rtimes C_4$.  Express it as a signed sum of letter contributions determined by prefix parity.  Which parts are covered by modular letter counts, and which require order-sensitive data?
\end{exercise}

\subsection{What a parser theorem must say}

A valid replacement for unique factorization may take several forms.
\begin{definition}[Canonical finite-state parser]
A parser is a deterministic transducer assigning to every word a sequence of marked factors such that the assignment is total, unique, and compositional enough that the desired modular statistic is a finite sum of factor contributions.
\end{definition}

\begin{definition}[Unambiguous rational decomposition]
A rational relation is unambiguous if each input word has at most one accepting output path.  If the relevant word language is the projection of an unambiguous relation with height-controlled fibers, one may avoid literal factor uniqueness.
\end{definition}

Any proposal must state how its output is turned into a generalized expression.  A parser that computes the group product is merely another automaton; the issue is height-one expression synthesis.

\section{Subgroups and actions of $A_5$}

\subsection{Concrete realization}

View $A_5$ as even permutations of five points.  It has order $60$ and conjugacy classes of sizes
\[
  1,\ 15,\ 20,\ 12,\ 12
\]
corresponding to identity, double transpositions, $3$-cycles, and the two split classes of $5$-cycles.  Mathlib formalizes this group and its simplicity \cite{mathlibAlternating2026}.

\subsection{Maximal subgroup types}

The maximal proper subgroups are of types
\[
  A_4\quad(|H|=12,\ [A_5:H]=5),
\]
\[
  D_5\quad(|H|=10,\ [A_5:H]=6),
\]
\[
  S_3\quad(|H|=6,\ [A_5:H]=10).
\]
They arise as normalizers/stabilizers of natural configurations.  Their coset actions give homomorphisms
\[
  A_5\to S_5,\qquad A_5\to S_6,\qquad A_5\to S_{10}.
\]
These actions can turn a group-recognized language into a permutation-automaton description.  Proper subgroups themselves are closer to known height-one families, except that $A_4$ is already the first unresolved barrier.

\subsection{Restriction and induction vocabulary}

For an $A_5$-module $V$ and subgroup $H$, restriction $\operatorname{Res}^G_HV$ forgets part of the action.  Induction $\operatorname{Ind}^G_HW$ builds a $G$-module from an $H$-module.  Mackey decomposition describes restriction of an induced module through double cosets.  These constructions may organize compatibility among local language data, but the language category is not automatically a module category.

The key question is not whether restriction and transfer exist in group cohomology.  It is whether there is an invariant of recognized languages on which these maps correspond to restriction and gluing of height-one certificates.

\begin{exercise}
For each maximal subgroup type, identify which Pin--Straubing--Th\'erien family covers it, if any.  Explain why coverage of all proper subgroups does not imply coverage of $A_5$.
\end{exercise}

\section{Representation theory as a source of finite coordinates}

\subsection{Permutation representations}

A finite $G$-set $X$ yields a permutation module $R[X]$.  For recognition, natural $G$-sets include automaton states, cosets $G/H$, conjugacy classes, and accepting fibers under translation.  A word path produces a sequence of basis vectors or transition operators.

Over characteristic zero, Maschke's theorem gives semisimplicity.  Over characteristics dividing $|G|$, extensions and cohomology become more visible.  Since generalized expressions are Boolean rather than linear objects, a representation-theoretic decomposition is only useful if one proves a map from Boolean language operations to the linear coordinates.

\subsection{Characters and class functions}

If accepting sets are unions of conjugacy classes, characters may compress their indicator functions.  But an arbitrary accepting subset need not be conjugacy invariant.  A theorem for class languages is a legitimate subclass result, not $\heightone(G)$.

A possible finite experiment is to decompose functions $G\to R$ under left-right translation.  Recognition fibers are basis vectors in this function space.  The question is whether height-one constructions correspond to a controlled subspace or filtration.  Complement and union are nonlinear over many coefficient rings, so the exact encoding matters.

\section{Group cohomology in low degrees}

\subsection{Cochains}

Let $G$ act on an abelian group $V$.  Inhomogeneous $n$-cochains are functions $G^n\to V$.  For a $1$-cochain $c$, the differential is
\[
  (dc)(g,h)=g\cdot c(h)-c(gh)+c(g).
\]
Thus $c$ is a cocycle exactly when
\[
  c(gh)=c(g)+g\cdot c(h).
\]
Coboundaries are $c(g)=g\cdot v-v$.  The quotient is $H^1(G,V)$.

For a $2$-cochain $f$, the cocycle identity is
\[
 g\cdot f(h,k)-f(gh,k)+f(g,hk)-f(g,h)=0.
\]
Factor sets of extensions satisfy this identity with the appropriate action.  The class in $H^2(Q,N)$, when $N$ is abelian, measures the obstruction to choosing a homomorphic section.

\subsection{Extensions and word evaluation}

Choose a set-theoretic section $s:Q\to G$.  Write $g=n s(q)$.  Then multiplication is
\[
(n,q)(n',q')=
\bigl(n\,(q\cdot n')\,f(q,q'),qq'\bigr).
\]
For a word, the $N$ component is a product of twisted letter terms and factor-set terms determined by adjacent prefix quotient states.  This gives a concrete path-count formula.  If $N$ is finite abelian, the terms may reduce to modular counts of transitions and pairs.

This is the strongest current reason to investigate cohomology: it organizes precisely the correction terms already appearing in extension normal forms.  It does not yet supply the generalized expression.

\subsection{Lyndon--Hochschild--Serre}

For $1\to N\to G\to Q\to1$, the LHS spectral sequence has
\[
  E_2^{p,q}=H^p\bigl(Q,H^q(N,V)\bigr)
  \Longrightarrow H^{p+q}(G,V)
\]
under standard hypotheses and conventions \cite{brown1982cohomology,weibel1994homological}.  Low-degree exact sequences relate invariants, restriction, inflation, and transgression.

A workshop should not begin by formalizing this spectral sequence.  First identify a coefficient module and a low-degree class that controls an explicit language construction.  Most candidate routes, if viable, should be testable in $H^1$ or $H^2$ before higher machinery.

\section{How to build a genuine language--cohomology bridge}

\subsection{Minimum specification}

A proposal must include:
\begin{enumerate}
  \item a finite alphabet $A$, recognizing morphism $\varphi:A^\ast\to G$, and accepting subset $P$;
  \item a $G$-module $V$ with an explicit action;
  \item a formula from words, prefix paths, or transition counts to a cochain;
  \item a proof of the cocycle identity;
  \item a theorem for at least one operation among union, complement, concatenation, and star;
  \item a statement of what vanishing or nonvanishing would imply about expressions.
\end{enumerate}

Without item 5, the invariant cannot be used to induct on height-one syntax.  Without item 6, it is not an obstruction or construction principle.

\subsection{Candidate 1: transition-count cocycles}

Let a quotient automaton have state set $Q$ and let $E=Q\times A$ be its directed labeled edges.  The free abelian group on $E$ records path counts.  Concatenating paths satisfies a twisted additive law because the second path begins at the endpoint of the first.  This naturally resembles a $1$-cocycle for a permutation action on edge coordinates.

Questions:
\begin{itemize}
  \item Does reduction modulo edge-cycle relations recover the normal component of a semidirect product?
  \item Can exact/modular edge-count fibers be expressed at height one?
  \item Does complement preserve the resulting class trivially, while concatenation requires a convolution law?
\end{itemize}

\subsection{Candidate 2: augmentation ideals}

For a group ring $R[G]$, the augmentation ideal $I_G$ is the kernel of the coefficient-sum map.  The element $\varphi(w)-1$ lies in $I_G$ and satisfies
\[
  gh-1=(g-1)+g(h-1),
\]
which is a cocycle identity for the left regular action.  This identity is exact and universal.

However, it only re-encodes group multiplication.  To affect star height, one must show that a filtration or cohomology class of $\varphi(w)-1$ has a controlled behavior under language constructors and separates some language.  The identity is a starting map, not the result.

\begin{exercise}[Augmentation cocycle]
Prove that $c(g)=g-1$ is a $1$-cocycle with values in the augmentation ideal under left multiplication.  Determine whether it is a coboundary.  Explain why this fact alone does not give a height-one expression.
\end{exercise}

\subsection{Candidate 3: syntactic bimodules}

Regular languages generally have noninvertible syntactic monoids.  A two-sided context acts naturally, suggesting bimodules and Hochschild cohomology rather than ordinary group cohomology.  This route may be more global, but it is also more distant from the current finite-group proofs.  A disciplined first target is a low-degree derivation attached to the syntactic morphism and a proof of its behavior under one closure operation.

\section{Pseudovarieties and wreath products}

Formal-language algebra often organizes classes of finite monoids as pseudovarieties, closed under finite products, submonoids, and quotients.  Wreath products encode cascades of automata and correspond to composition principles.  Pin--Straubing--Th\'erien obtain height-one results for pseudovariety constructions involving aperiodic monoids and commutative groups \cite{pinStraubingTherien1992}.

For group theorists, a wreath product resembles an induced semidirect product with a function group.  For formalization, the main danger is abstracting too early.  First formalize the exact certificate transport for the specific product used in the proof.  Then package it as a pseudovariety closure theorem.

\begin{exercise}
Take a cascade of a star-free automaton and a cyclic group automaton.  Write the transition monoid as a divisor of a wreath product and state the height-one transport theorem needed.  Identify every closure operation invoked.
\end{exercise}

\section{Lower bounds from an algebraist's perspective}

A disproof needs an invariant of languages, not merely a complicated group.  Algebraic candidates might include profinite identities, derived categories, cohomological dimensions, or a filtration on syntactic algebras.  The required theorem is closure under all height-one constructors.

A useful test protocol is:
\begin{enumerate}
  \item compute the invariant for star-free languages;
  \item prove its behavior under one star of a star-free operand;
  \item prove closure under Boolean operations;
  \item attack concatenation, the likely failure point;
  \item test against cyclic, abelian, class-two, and $A\rtimes C_2$ known families;
  \item only then test $A_4$ or $A_5$ candidates.
\end{enumerate}
A counterexample to the proposed closure law should be preserved; it rules out an entire lower-bound mechanism.

\section{Lean-facing formulation of group facts}

\subsection{Concrete versus abstract models}

Mathlib already represents $A_n$ as a subgroup of permutations and includes finite and simplicity results for $A_5$ \cite{mathlibAlternating2026}.  Use this model unless a proof requires a semidirect-product model.  If two models are used, prove an explicit multiplicative equivalence and transport the recognition property.

For $\operatorname{Dic}_3$, agree on a canonical model before proving language theorems.  A presentation alone is not enough unless the order and universal property are established.  A semidirect product may be algebraically convenient but can carry API overhead; a permutation subgroup may be easier for finite computation.

\subsection{Theorem signatures}

A useful group theorem should expose the data it transforms.  Schematic examples are:
\begin{lstlisting}[language={},caption={Lean-facing algebraic interfaces.}]
def HeightOneForGroup (G : Type v) [Group G] : Prop :=
  forall (A : Type u) [Fintype A] [DecidableEq A],
    forall R : Recognition A G,
      HasHeightAtMost R.language 1

structure ExtensionLanguageData where
  phi : List A ->* G
  accept : Set G
  normal : Subgroup G
  normal_is_normal : normal.Normal
  quotient_certificate : ...
  cocycle_certificate : ...
\end{lstlisting}
The second structure should only be introduced after a concrete proof determines its fields.  Do not design an abstract extension API from analogy alone.

\section{Workshop protocols}

\subsection{For a proposed group-theoretic lemma}

The proposer writes:
\begin{enumerate}
  \item exact group hypotheses;
  \item exact word-evaluation formula;
  \item the finite statistics required;
  \item the language operation realizing each statistic;
  \item the generalized-expression height calculation;
  \item smallest nontrivial and degenerate examples;
  \item a falsification test for factorization or gluing.
\end{enumerate}
A formal-language theorist then restates the claim without group shorthand.  A Lean expert translates the quantifiers.  Disagreement is resolved before proof search.

\subsection{For an AI-generated cohomology idea}

Reject immediately if it lacks coefficients, action, and a word-to-cochain map.  Otherwise:
\begin{enumerate}
  \item verify the cocycle equation by exact finite computation;
  \item test all language constructors on the smallest examples;
  \item ask whether the missing theorem is equivalent to height-one collapse;
  \item preserve the first failing example;
  \item only after survival, consult deeper cohomological machinery.
\end{enumerate}

\section{Exercises with workshop outputs}

\begin{exercise}[Semidirect word formula]
For $N\rtimes Q$, derive the normal component of a word in terms of prefix products in $Q$.  Package it as a path statistic on the Cayley automaton of $Q$.
\end{exercise}

\begin{exercise}[Class-two language]
Choose a class-two finite group and an accepting fiber.  Produce the finite list of modular count conditions that characterize the language.  Hand this list to the certificate team.
\end{exercise}

\begin{exercise}[$C_3$ ambiguity]
Formulate the factorization used in the attempted $A\rtimes C_3$ extension.  Search for the smallest word with two valid decompositions and explain how the resulting counts differ.
\end{exercise}

\begin{exercise}[$A_4$ exhaustive reduction]
Classify generating pairs of $A_4$ up to automorphism.  State precisely what further symmetries act on accepting subsets.  Do not claim exhaustive language coverage until the reduction is proved.
\end{exercise}

\begin{exercise}[$A_5$ coset actions]
Construct the actions on cosets of $A_4$, $D_5$, and $S_3$.  For a fixed two-letter morphism, compute the corresponding permutation automata and compare their information content.
\end{exercise}

\begin{exercise}[Factor set]
For a nonsplit or non-homomorphic section of a small extension with abelian kernel, compute the factor set and verify the $2$-cocycle identity.  Translate its terms into quotient-path transitions.
\end{exercise}

\begin{exercise}[Invariant attack]
Propose a group-cohomological quantity for recognized languages.  Test union, complement, concatenation, and star separately.  The deliverable may be a counterexample showing the quantity cannot be a height-one invariant.
\end{exercise}

\section{Reading guide}

For the formal-language algebra, read Pin--Straubing--Th\'erien \cite{pinStraubingTherien1992}, then Bourne--Ru\v{s}kuc and Bourne's thesis \cite{bourneRuskuc2016,bourne2017counting}.  For finite-group background, a standard group text such as Rotman is sufficient \cite{rotman1995introduction}.  Brown is the primary cohomology reference for this workshop, with Weibel for homological organization \cite{brown1982cohomology,weibel1994homological}.  Consult these only in service of a specified language map.

\subsection*{Acknowledgement}

This primer uses the package, theorem-environment, bibliography, and macro organization of the supplied Ryuya Hora \code{FirstVersion.tex} template, with portable-font and workshop-specific additions.

\printbibliography
\end{document}
